% Insert inside the manuscript's Discussion section.
The network reveals a reproducible association between ingredient similarity and cuisine tags, together with nonuniform links between cuisines. The ingredient-level randomization narrows this interpretation: most of the large within-cuisine enrichment relative to label shuffling is consistent with cuisine-specific ingredient frequencies, but the observed network retains an additional co-occurrence signal. This residual persists after excluding recipes explicitly named as sweets or drinks and after removing very high-similarity edges. It also persists when ingredient exchanges are restricted to recipes sharing a broad Food.com dish type. The Chinese--Japanese pair illustrates why the reference model matters: its strong relative affinity under label shuffling disappears when ingredient frequencies within cuisines are preserved. Greek--Middle Eastern links retain an exploratory excess under the stricter comparisons. These edges should be interpreted as shared ingredient lists, not as evidence of culinary transmission or recipe authenticity. Food.com cuisine tags are user supplied and can be inconsistent: one Chinese-tagged recipe in the manual review was explicitly titled ``Benihana Japanese Fried Rice,'' while a Greek-tagged hummus recipe described itself as Lebanese. Generic batters and near-duplicate ingredient lists also produced high-similarity links. The qualitative examples therefore illustrate possible mechanisms behind the aggregate pattern rather than validating every edge. Network estimates depend on ingredient normalization and sparsification choices; the reported threshold checks support the direction of the within-cuisine result but do not remove selection or labeling bias. Food.com course tags can overlap and be ambiguous, and the type-stratified null still disrupts finer preparation styles and near-duplicate recipe families. Its residual cannot be attributed uniquely to cuisine.
